\documentclass[10pt]{article}
\usepackage[T1]{fontenc}
\usepackage[utf8]{inputenc}
\usepackage[french]{babel}
\usepackage{amsmath,amssymb,mathrsfs,mathpazo}
\usepackage{geometry}
\geometry{a4paper,margin=2cm}
\usepackage{graphicx,array,tabularx,tabularray,colortbl,multirow,multicol,booktabs,enumitem,pifont,ulem,fancybox,diagbox,cancel}
\usepackage{pstricks,pstricks-add,pst-plot,pst-tree,pst-node,pst-text,pst-func,pst-math,pst-bspline,pst-3dplot,pst-eucl,pst-solides3d}
\let\pscalculate\undefined
\usepackage{pst-bezier}
\usepackage{tikz,pgfplots,tkz-tab}
\usetikzlibrary{arrows,arrows.meta,patterns,calc,positioning,trees,decorations.markings,shapes.geometric,shapes.misc,shapes.arrows,intersections,decorations.pathmorphing,decorations.pathreplacing,backgrounds,shadows,babel}
\pgfplotsset{compat=1.18}
\usepgfplotslibrary{fillbetween}
\usepackage[np]{numprint}
\usepackage{listings,makecell,pas-tableur,twoopt}
\let\barre\undefined
\usepackage[pstricks]{tablvar}
\definecolor{darkgreen}{RGB}{0,100,0}
\providecommand{\red}{\color{red}}
\providecommand{\blue}{\color{blue}}
\providecommand{\green}{\color{green}}
\providecommand{\black}{\color{black}}
\providecommand{\white}{\color{white}}
\providecommand{\gray}{\color{gray}}
\providecommand{\cyan}{\color{cyan}}
\providecommand{\magenta}{\color{magenta}}
\providecommand{\yellow}{\color{yellow}}
\providecommand{\violet}{\color{violet}}
\providecommand{\orange}{\color{orange}}
\providecommand{\euro}{\texteuro}
\providecommand{\EUR}{\texteuro}
\providecommand{\ieme}{\textsuperscript{e}}
\providecommand{\degres}{\ensuremath{^\circ}}
\providecommand{\degree}{\ensuremath{^\circ}}
\providecommand{\vv}[1]{\overrightarrow{#1}}
\providecommand{\Coord}[1]{#1}
\providecommand{\decosix}{}
\providecommand{\starredbullet}{\ensuremath{\bullet}}
\pagestyle{empty}
\setlength{\parindent}{0pt}
\usepackage{adjustbox}\DeclareUnicodeCharacter{2236}{:}\DeclareUnicodeCharacter{2190}{\ensuremath{\leftarrow}}
\providecommand{\R}{\mathbb{R}}
\providecommand{\N}{\mathbb{N}}
\providecommand{\Z}{\mathbb{Z}}
\providecommand{\Q}{\mathbb{Q}}
\providecommand{\D}{\mathbb{D}}
\providecommand{\C}{\mathbb{C}}
\providecommand{\e}{{\mathrm{e}}}
\providecommand{\vect}[1]{\overrightarrow{#1}}
\providecommand{\vectt}[1]{\overrightarrow{\mathrm{#1}}}
\providecommand{\vvt}[1]{\overrightarrow{\mathrm{#1}}}
\providecommand{\np}[2][]{{#2~#1}}
\providecommand{\vv}[1]{{\overrightarrow{#1}}}
\renewcommand{\d}{{\mathrm{d}}}
\providecommand{\degres}{{^\circ}}
\providecommand{\degre}{{^\circ}}
\providecommand{\blue}{\color{blue}}
\providecommand{\red}{\color{red}}
\providecommand{\enskip}{\quad}
\providecommand{\llbracket}{\lbrack\!\lbrack}
\providecommand{\rrbracket}{\rbrack\!\rbrack}
\providecommand{\euro}{{\text{€}}}
\providecommand{\no}{{\text{n°}}}
\providecommand{\interval}[2]{{[#1;#2]}}
\providecommand{\intervFO}[2]{{[#1;#2[}}
\providecommand{\dotfill}{\ldots}

\providecommand{\ds}{\displaystyle}
\providecommand{\textding}[1]{#1}
\providecommand{\Oij}{\(\mathrm{O};\vec{\imath},\vec{\jmath}\)}
\providecommand{\Oijk}{\(\mathrm{O};\vec{\imath},\vec{\jmath},\vec{k}\)}
\providecommand{\Ouv}{\(\mathrm{O};\vec{u},\vec{v}\)}
\providecommand{\intervFF}[2]{\ensuremath{[#1;#2]}}
\renewcommand{\euro}{\ifmmode\text{\texteuro}\else\texteuro\fi}
\DeclareUnicodeCharacter{00B0}{\ifmmode^{\circ}\else\textdegree\fi}
\DeclareUnicodeCharacter{20AC}{\ifmmode\text{\texteuro}\else\texteuro\fi}
% Énoncés : Éducation nationale. Transcription et corrigés : APMEP, Denis Vergès et contributeurs. Découpage et classement : CoopMaths, Sébastien Lozano, Éric Elter et contributeurs. Adaptation de la présentation : Maths Explorer.
% CC BY-SA : https://coopmaths.fr/www/about/
% https://coopmaths.fr/alea/static/e3c/2020/tex/e3c_2020_00_sujet18_2.tex
\begin{document}
\noindent\textit{APMEP / CoopMaths. CC BY-SA. Présentation adaptée par Maths Explorer.}\par\medskip
\medskip

Une banque propose un placement. Le compte est rémunéré et rapporte 5\,\% par an. La banque
prend des frais de gestion qui se montent à 12 euros par an.

Ainsi, chaque année la somme sur le compte augmente de 5\,\% puis la banque prélève 12 euros.

Noémie place la somme de \np{1000} euros dans cette banque.

On appelle $u_n$ la somme disponible sur le compte en banque de Noémie après $n$ années, où
$n$ désigne un entier naturel.

On a donc $u_0 = \np{1000}$ et pour tout entier naturel n, $u_{n+1} = 1,05 u_n-12$

\begin{enumerate}
\item  Avec un tableur on a calculé les premiers termes de la suite
$\left(u_n\right)$ :

\medskip

\begin{minipage}[]{0.425\textwidth}
\begin{enumerate}
\item Quelle formule a-t-on entrée dans la cellule B3 avant de
l’étirer pour obtenir ces résultats ?
\item En utilisant les valeurs calculées de la suite, indiquer à
Noémie combien de temps elle doit attendre pour que son
placement lui rapporte 20\,\%.
\end{enumerate}
\end{minipage}
\hspace{2cm}
\begin{minipage}[]{8.97cm}
\begin{tabularx}{0.445\textwidth}{|>{\columncolor{lightgray}}c|*{2}{>{\raggedleft \arraybackslash}X|}}
\hline
\rowcolor{lightgray}&\centering  A& \centering B\tabularnewline\hline
1&\raggedright n&\raggedright u(n)\tabularnewline\hline
2&0&\np{1000}\\\hline
3&1&\np{1038.00}\\\hline
4&2&\np{1077.90}\\\hline
5&3&\np{1119.80}\\\hline
6&4&\np{1163.78}\\\hline
7&5&\np{1209.97}\\\hline
8&6&\np{1258.47}\\\hline
9&7&\np{1309.40}\\\hline
10&8&\np{1362.87}\\\hline
11&9&\np{1419.01}\\\hline
12&10&\np{1477.96}\\\hline
\end{tabularx}
\end{minipage}

\end{enumerate}

\medskip

On pose $\left(v_n\right)$ la suite définie, pour tout entier naturel $n$, par $v_n = u_n- 240$.

\begin{enumerate}[resume]
\item Montrer que la suite $\left(v_n\right)$ est géométrique de raison $1,05$.
\item Exprimer $v_n$ puis $u_n$ en fonction de l’entier $n$.
\item Calculer à partir de cette dernière formule la somme disponible sur le compte en banque de
Noémie après 20 ans de placement.
\end{enumerate}

\vspace{0,5cm}

\end{document}
